\begin{solapartat}
\punts{0,1} Com que el nombre de protons és 92 i el nombre nucleons és 235 llavors el nombre de neutrons és 143.

\punts{0,4} El defecte de massa és:
\[ \Delta m = \bigl(92\,m({}^{1}_{1}p) + 143\,m({}^{1}_{0}n)\bigr) - m({}^{235}_{92}U) = 3,18\times 10^{-27}\ \mathrm{kg} \]

\punts{0,3} I l'energia d'enllaç és
\begin{align*}
E = \Delta mc^2 &= 3,18\times 10^{-27}\times (3,00\times 10^{8})^2\\
&= 2,862\times 10^{-10}\ \mathrm{J} = 1,78\times 10^{9}\ \mathrm{eV}
\end{align*}

\punts{0,2} I l'energia d'enllaç per nucleó:
\[ \frac{\Delta mc^2}{A} = \frac{2,862\times 10^{-10}\ \mathrm{J}}{235} = 1,22\times 10^{-12}\ \mathrm{J} = 7,60\times 10^{6}\ \mathrm{eV} \]

\punts{0,25} A partir de la taula:

\begin{center}
\begin{tabular}{|l|l|l|l|l|}
\hline
Nucli & Sofre-34 & \textcolor{red}{\textbf{Ferro-56}} & Radi-226 & Urani-235 \\
\hline
Energia d'enllaç per nucleó (MeV) & 8,58 & \textcolor{red}{\textbf{8,79}} & 7,66 & \destaca{7,60} \\
\hline
\end{tabular}
\end{center}

Com que l'estabilitat dels nuclis augmenta amb l'energia d'enllaç per nucleó, podem comprovar que el nucli més estable de la taula seria el Ferro-56.
\end{solapartat}

\begin{solapartat}
\punts{0,6} Reacció nuclear
\[ {}^{235}_{92}U + {}^{1}_{0}n \rightarrow {}^{141}_{56}Ba + 3\,{}^{1}_{0}n + {}^{x}_{y}Z \]
\[ \left.\begin{aligned}
235 + 1 &= 141 + 3 + x \longrightarrow x = 92\\
92 &= 56 + y \longrightarrow y = 36
\end{aligned}\right\}\ {}^{92}_{36}Kr \]
Per tant, la reacció és
\[ {}^{235}_{92}U + {}^{1}_{0}n \rightarrow {}^{141}_{56}Ba + {}^{92}_{36}Kr + 3\,{}^{1}_{0}n \]

\punts{0,15} Energia consumida per la il·luminació d'un estadi esportiu en MeV:
\[ 25\,000\ \mathrm{kWh} = 25\,000\times 1000\times 3600 = 9,00\times 10^{10}\ \mathrm{J} \]
\[ 9,00\times 10^{10}\ \mathrm{J}\,\frac{1\ \mathrm{eV}}{1,602\times 10^{-19}\ \mathrm{J}}\,\frac{1\ \mathrm{MeV}}{10^{6}\ \mathrm{eV}} = 5,62\times 10^{23}\ \mathrm{MeV} \]

\punts{0,25} Nuclis d'urani
\[ n({}^{235}_{92}U) = \frac{5,62\times 10^{23}}{202,5} = 2,77\times 10^{21} \]

\punts{0,25} Massa d'urani
\[ n({}^{235}_{92}U)\,\frac{3,902158\times 10^{-25}\ \mathrm{kg}}{1\ \text{nucli}\ {}^{235}_{92}U} = 1,083\times 10^{-3}\ \mathrm{kg} = 1,083\ \mathrm{g} \]
\end{solapartat}
